% sections/conclusion.tex — Conclusion
\section{Conclusion}
\label{sec:conclusion}

We have presented \system{}, a framework that addresses a fundamental
challenge in enterprise audit analytics: the absence of ground truth.
We showed that constructing knowledge systems primarily from operational enterprise
data---while valuable and widely practiced---faces inherent
limitations when the source data itself may contain systematic errors,
leaving teams without a comprehensive reference against which to
verify the resulting knowledge structures.  When the source data contains systematic errors, the
constructed knowledge may inherit those errors, potentially masking the
root causes that audit procedures seek to identify.

Our formal analysis demonstrated three results:
\textbf{(i)}~recovering the true accounting network from observed
journal entries is combinatorially infeasible, with the configuration
space growing super-exponentially (\Cref{thm:combinatorial});
\textbf{(ii)}~systematic errors in source data propagate through
multi-stage enterprise processes with high probability of remaining
undetected (\Cref{prop:error_persistence}); and
\textbf{(iii)}~the inter-dimensional nature of enterprise data---spanning
topology, statistics, temporality, norms, and organization---means that
errors in one dimension amplify across others
(\Cref{thm:amplification}).

\system{} addresses this through \emph{forward generation} from a
fully specified model, producing \emph{reference knowledge graphs}
where the complete audit trail is known by construction across three
layers of knowledge: structural (entity types and relationship types
covering major enterprise processes), statistical (Benford-compliant
distributions, copula-based dependencies), and normative (multiple
accounting frameworks, ISA/PCAOB/SOX audit standards, COSO~2013
internal controls).

Experimental evaluation confirms that generated datasets satisfy all
hard accounting invariants at 100\% compliance, achieve Benford MAD
scores below 0.006, and support downstream fraud detection models
achieving F1 scores of 0.84 with GNNs.  The system provides full provenance for every record in the reference
knowledge graph, with generation throughput and domain coverage detailed
in \Cref{sec:experiments}.

% ─────────────────────────────────────────────────────────────────────
\subsection{Limitations}
\label{sec:limitations}

\begin{itemize}
  \item \textbf{Domain coverage.}  While \system{} covers 20+ enterprise
    processes, specialized domains such as insurance claims, healthcare
    billing, and government accounting are not yet modeled.

  \item \textbf{Learned distributions.}  The current parametric approach
    (log-normal mixtures, copulas) could be complemented by deep
    generative models for more complex distributional structures, at
    the cost of reproducibility and interpretability.

  \item \textbf{Cross-enterprise patterns.}  The fingerprint module
    operates on single enterprises.  Cross-organization patterns
    (industry benchmarking, market-level correlations) are not yet
    modeled.

  \item \textbf{Validation against real knowledge graphs.}  While
    \system{}'s statistical properties are calibrated against 155
    real-world datasets, the reference knowledge graph structure
    has not been validated against actual enterprise knowledge graphs
    (which are themselves subject to the ground truth problem).
\end{itemize}

% ─────────────────────────────────────────────────────────────────────
\subsection{Future Work}
\label{sec:future_work}

Several of the research directions identified in earlier versions of
this paper have been realized in v2.0.0: the audit FSM engine now
includes 12 blueprints spanning ISA, IIA-GIAS, PCAOB, AICPA SOC~2,
and Big~4 firm-style methodologies; the audit-optimizer crate performs
shortest-path analysis and Monte Carlo simulation over blueprint
graphs; and conformance metrics are computed across methodology
variants.  The following directions remain open:

\begin{enumerate}
  \item \textbf{Temporal graph network export.}  Integrating temporal
    graph network (TGN) export and the temporal accounting network
    formalism of~\citet{ivertowski2024temporal} for time-evolving
    reference knowledge graphs with continuous-time event encoding.

  \item \textbf{Causal inference from synthetic ground truth.}  Using
    \system{}'s fully provenanced datasets---where the causal DAG is
    known by construction---as benchmarks for evaluating causal
    discovery algorithms in the enterprise domain.

  \item \textbf{Federated simulation.}  Extending the federated
    fingerprint aggregation framework to support distributed
    multi-organization simulation under secure multi-party computation,
    enabling industry-level reference knowledge graphs without
    centralizing enterprise data.

  \item \textbf{Hardware-accelerated graph construction.}  Coupling
    \system{}'s output with GPU-accelerated accounting network
    generation~\citep{ivertowski2024hardware} for interactive audit
    analytics.

  \item \textbf{Adaptive anomaly calibration.}  Using reinforcement
    learning to tune anomaly injection such that downstream detector
    performance matches a target difficulty curve.

  \item \textbf{Knowledge graph completion benchmarks.}  Using
    \system{}'s reference graphs to create standardized benchmarks for
    evaluating knowledge graph construction and completion algorithms
    in the audit domain.

  \item \textbf{Learned causal transfer functions.}  Replacing
    parametric transfer functions with neural network-based functions
    fitted from real-world data for more faithful counterfactual
    propagation.

  \item \textbf{Judgment automation boundary analysis.}  Leveraging
    the \texttt{judgment\_level} annotations across 12 methodology
    blueprints to empirically characterize the automation frontier in
    audit procedures---identifying which steps are amenable to AI
    assistance and where human judgment remains indispensable.
\end{enumerate}

% ─────────────────────────────────────────────────────────────────────
\subsection{Availability}
\label{sec:availability}

\system{} is released as open-source software under the Apache~2.0
license.  The source code, documentation, example configurations,
and pre-built binaries are available at
\url{https://github.com/mivertowski/SyntheticData}.
A Python wrapper (\texttt{datasynth-py}) provides high-level access
via \texttt{pip install datasynth-py}.
